\begin{thebibliography}{99}\small
\bibitem{R01} Subramanian A;  Narayan R;  Corsello SM; et al.. A Next Generation Connectivity Map: L1000 Platform and the First 1,000,000 Profiles. Cell, 2017. DOI: \nolinkurl{10.1016/j.cell.2017.10.049}.
\bibitem{R02} Lambert SA;  Jolma A;  Campitelli LF; et al.. The Human Transcription Factors. Cell, 2018. DOI: \nolinkurl{10.1016/j.cell.2018.01.029}.
\bibitem{R03} Bianchi E;  Zini R;  Salati S; et al.. c-myb supports erythropoiesis through the transactivation of KLF1 and LMO2 expression. Blood, 2010. DOI: \nolinkurl{10.1182/blood-2009-08-238311}.
\bibitem{R04} Vegiopoulos A;  García P;  Emambokus N; et al.. Coordination of erythropoiesis by the transcription factor c-Myb. Blood, 2006. DOI: \nolinkurl{10.1182/blood-2005-07-2968}.
\bibitem{R05} Tunyaplin C;  Shaffer AL;  Angelin-Duclos CD; et al.. Direct repression of prdm1 by Bcl-6 inhibits plasmacytic differentiation. Journal of immunology (Baltimore, Md. : 1950), 2004. DOI: \nolinkurl{10.4049/jimmunol.173.2.1158}.
\bibitem{R06} Date D;  Das R;  Narla G; et al.. Kruppel-like transcription factor 6 regulates inflammatory macrophage polarization. The Journal of biological chemistry, 2014. DOI: \nolinkurl{10.1074/jbc.M113.526749}.
\bibitem{R07} Kim GD;  Das R;  Goduni L; et al.. Kruppel-like Factor 6 Promotes Macrophage-mediated Inflammation by Suppressing B Cell Leukemia/Lymphoma 6 Expression. The Journal of biological chemistry, 2016. DOI: \nolinkurl{10.1074/jbc.M116.738617}.
\bibitem{R08} Lin H;  Grosschedl R. Failure of B-cell differentiation in mice lacking the transcription factor EBF. Nature, 1995. DOI: \nolinkurl{10.1038/376263a0}.
\bibitem{R09} Frangoul H;  Altshuler D;  Cappellini MD; et al.. CRISPR-Cas9 Gene Editing for Sickle Cell Disease and $\beta$-Thalassemia. The New England journal of medicine, 2021. DOI: \nolinkurl{10.1056/NEJMoa2031054}.
\bibitem{R10} Breda L;  Papp TE;  Triebwasser MP; et al.. In vivo hematopoietic stem cell modification by mRNA delivery. Science (New York, N.Y.), 2023. DOI: \nolinkurl{10.1126/science.ade6967}.
\bibitem{R11} Norman TM;  Horlbeck MA;  Replogle JM; et al.. Exploring genetic interaction manifolds constructed from rich single-cell phenotypes. Science, 2019. DOI: \nolinkurl{10.1126/science.aax4438}.
\bibitem{R12} Roohani Y;  Huang K;  Leskovec J. Predicting transcriptional outcomes of novel multigene perturbations with GEARS. Nature biotechnology, 2024. DOI: \nolinkurl{10.1038/s41587-023-01905-6}.
\bibitem{R13} Cui H;  Wang C;  Maan H; et al.. scGPT: toward building a foundation model for single-cell multi-omics using generative AI. Nature methods, 2024. DOI: \nolinkurl{10.1038/s41592-024-02201-0}.
\bibitem{R14} Hao M;  Gong J;  Zeng X; et al.. Large-scale foundation model on single-cell transcriptomics. Nature methods, 2024. DOI: \nolinkurl{10.1038/s41592-024-02305-7}.
\bibitem{R15} Ahlmann-Eltze C;  Huber W;  Anders S. Deep-learning-based gene perturbation effect prediction does not yet outperform simple linear baselines. Nature Methods, 2025. DOI: \nolinkurl{10.1038/s41592-025-02772-6}.
\bibitem{R16} Wenkel F;  Tu W;  Masschelein C; et al.. TxPert: using multiple knowledge graphs for prediction of transcriptomic perturbation effects. Nature Biotechnology, 2026. DOI: \nolinkurl{10.1038/s41587-026-03113-4}.
\end{thebibliography}